% GENERATED FILE. Edit canonical lessons, then run npm run generate:books.
% canonical-source-hash: fnv1a64:0e1ec829b830610b
% canonical-lessons: MR-C132-hatoda, MR-C132-katri, MR-C132-sui, MR-C132-dora, MR-C132-kangva

\chapter{Hammer, Scissors, Needle, Thread, Comb}
\label{ch:mr-hammer-scissors-needle-thread-comb}
\hlchaptermodality{132}

\begin{chapteropening}
\textbf{By the end of this chapter:} \emph{I can say a hammer, scissors, a needle, a thread and a comb.}

Complete the last lesson of chapter 132: i can say a hammer, scissors, a needle, a thread and a comb.
\end{chapteropening}

\section[hātoḍā]{\mr{हातोडा} (hātoḍā) — a hammer}

\label{lesson:MR-C132-hatoda}



\begin{quote}
Before the new one: say the Marathi for a rope, then the Marathi for a nail.
\end{quote}

\subsection*{You'll want to know: \mr{हातोडा}}
\textbf{\mr{हातोडा}} — \emph{hātoḍā} — \textquotedblleft{}a hammer\textquotedblright{}.

\begin{grammarlens}[title={What you've built}]
One more word for this chapter.
\end{grammarlens}

\subsection*{Your turn}
\begin{itemize}
\raggedright
  \item \emph{Say it:} \emph{hātoḍā}
  \item \emph{Say it:} \emph{hātoḍā}, once more
  \item \emph{Say it:} say \emph{hātoḍā}
  \item \emph{Recall:} say the Marathi for a rope, then the Marathi for a nail, then say \emph{hātoḍā} again
\end{itemize}

\begin{tcolorbox}[breakable,colback=teal!4,colframe=teal!35!black,title={Before you move on}]
What does \mr{हातोडा} mean? (\textquotedblleft{}A hammer\textquotedblright{}.) Say it once more.
\end{tcolorbox}
\section[kātrī]{\mr{कात्री} (kātrī) — scissors}

\label{lesson:MR-C132-katri}



\begin{quote}
Before the new one: say the Marathi for a nail, then the Marathi for a hammer.
\end{quote}

\subsection*{You'll want to know: \mr{कात्री}}
\textbf{\mr{कात्री}} — \emph{kātrī} — \textquotedblleft{}scissors\textquotedblright{}.

\begin{grammarlens}[title={What you've built}]
One more word for this chapter.
\end{grammarlens}

\subsection*{Your turn}
\begin{itemize}
\raggedright
  \item \emph{Say it:} \emph{kātrī}
  \item \emph{Say it:} \emph{kātrī}, once more
  \item \emph{Say it:} say \emph{kātrī}
  \item \emph{Recall:} say the Marathi for a nail, then the Marathi for a hammer, then say \emph{kātrī} again
\end{itemize}

\begin{tcolorbox}[breakable,colback=teal!4,colframe=teal!35!black,title={Before you move on}]
What does \mr{कात्री} mean? (\textquotedblleft{}Scissors\textquotedblright{}.) Say it once more.
\end{tcolorbox}
\section[suī]{\mr{सुई} (suī) — a needle}

\label{lesson:MR-C132-sui}



\begin{quote}
Before the new one: say the Marathi for a hammer, then the Marathi for scissors.
\end{quote}

\subsection*{You'll want to know: \mr{सुई}}
\textbf{\mr{सुई}} — \emph{suī} — \textquotedblleft{}a needle\textquotedblright{}.

\begin{grammarlens}[title={What you've built}]
One more word for this chapter.
\end{grammarlens}

\subsection*{Your turn}
\begin{itemize}
\raggedright
  \item \emph{Say it:} \emph{suī}
  \item \emph{Say it:} \emph{suī}, once more
  \item \emph{Say it:} say \emph{suī}
  \item \emph{Recall:} say the Marathi for a hammer, then the Marathi for scissors, then say \emph{suī} again
\end{itemize}

\begin{tcolorbox}[breakable,colback=teal!4,colframe=teal!35!black,title={Before you move on}]
What does \mr{सुई} mean? (\textquotedblleft{}A needle\textquotedblright{}.) Say it once more.
\end{tcolorbox}
\section[dorā]{\mr{दोरा} (dorā) — a thread}

\label{lesson:MR-C132-dora}



\begin{quote}
Before the new one: say the Marathi for scissors, then the Marathi for a needle.
\end{quote}

\subsection*{You'll want to know: \mr{दोरा}}
\textbf{\mr{दोरा}} — \emph{dorā} — \textquotedblleft{}a thread\textquotedblright{}.

\begin{grammarlens}[title={What you've built}]
One more word for this chapter.
\end{grammarlens}

\subsection*{Your turn}
\begin{itemize}
\raggedright
  \item \emph{Say it:} \emph{dorā}
  \item \emph{Say it:} \emph{dorā}, once more
  \item \emph{Say it:} say \emph{dorā}
  \item \emph{Recall:} say the Marathi for scissors, then the Marathi for a needle, then say \emph{dorā} again
\end{itemize}

\begin{tcolorbox}[breakable,colback=teal!4,colframe=teal!35!black,title={Before you move on}]
What does \mr{दोरा} mean? (\textquotedblleft{}A thread\textquotedblright{}.) Say it once more.
\end{tcolorbox}
\section[kaṅgvā]{\mr{कंगवा} (kaṅgvā) — a comb}

\label{lesson:MR-C132-kangva}



\begin{quote}
Before the new one: say the Marathi for a needle, then the Marathi for a thread.
\end{quote}

\subsection*{You'll want to know: \mr{कंगवा}}
\textbf{\mr{कंगवा}} — \emph{kaṅgvā} — \textquotedblleft{}a comb\textquotedblright{}.

\begin{grammarlens}[title={What you've built}]
One more word for this chapter.
\end{grammarlens}

\subsection*{Your turn}
\begin{itemize}
\raggedright
  \item \emph{Say it:} \emph{kaṅgvā}
  \item \emph{Say it:} \emph{kaṅgvā}, once more
  \item \emph{Say it:} say \emph{kaṅgvā}
  \item \emph{Recall:} say the Marathi for a needle, then the Marathi for a thread, then say \emph{kaṅgvā} again
\end{itemize}

\begin{tcolorbox}[breakable,colback=teal!4,colframe=teal!35!black,title={Before you move on}]
What does \mr{कंगवा} mean? (\textquotedblleft{}A comb\textquotedblright{}.) Say it once more.
\end{tcolorbox}
